\documentclass[12pt,a4paper]{article}
\usepackage{ugmath}
\usepackage{float}
\usepackage{placeins}
\usepackage{array}
\usepackage[labelfont=bf,labelsep=space]{caption}
\newcommand{\paren}[1]{\left( #1 \right)}
\newcommand{\alumno}{Ricardo León Martínez}
\newcommand{\materia}{Elementos de Probabilidad y Estadistica}
\newcommand{\profesor}{Claudia Reynoso Alcántara}
\newcommand{\tarea}{Tarea 5}
\newcommand{\fecha}{25/3/2026}

\begin{document}

    \begin{center}
        {\large\textbf{UNIVERSIDAD DE GUANAJUATO}}\\[0.3cm]
        {\normalsize\textbf{DIVISIÓN DE CIENCIAS NATURALES Y EXACTAS}}\\
        {\normalsize\textbf{CAMPUS GUANAJUATO}}\\[1cm]

        {\Large\textbf{\tarea\ (\materia)}}\\[1cm]
    \end{center}

    \noindent
    \textbf{Nombre:} \alumno 
    \hfill 
    \textbf{Fecha:} \fecha 
    \hfill 
    \textbf{Calificación:} \rule{3cm}{0.4pt}

    \vspace{0.3cm}

    \begin{mdframed}[style=mdpinkbox,frametitle={Ejercicio 1}]
        Pedro tiene una casa a la cual considera ponerle una alarma antirrobos. Sea $A$ el evento de
        que Pedro instale una alarma antes de finalizar el año y sea $B$ el evento de que roben la casa
        de Pedro antes de que finalice el año. Conteste lo siguiente.

        (a) Intuitivamente, ¿cuál probablidad entre $P(A \mid B)$ o $P(A \mid B^c)$ es más grande? Explique su respuesta.

        (b) Intuitivamente, ¿cuál probablidad entre $P(B \mid A)$ o $P(B \mid A^c)$ es más grande? Explique su respuesta.

        (c) Demuestre que para cualesquiera eventos $A, B$ con probabilidades $0 < P(A), P(B) < 1$, se cumple que
        \[
        P(A \mid B) > P(A \mid B^c) \iff P(B \mid A) > P(B \mid A^c).
        \]

        (d) En 1982, Tversky y Kahneman publicaron una encuesta en la que se le preguntó los incisos (a) y (b)
        a 162 personas, de las cuales 131 respondieron que $P(A \mid B) > P(A \mid B^c)$ pero $P(B \mid A) < P(B \mid A^c)$.
        Dé una posible explicación del porqué esta respuesta fue tan popular a pesar de que estas dos
        desigualdades no pueden satisfacerse simultáneamente.
    \end{mdframed}

    \begin{mdframed}[style=mdpinkbox,frametitle={Ejercicio 2}]
        Decimos que $B$ aporta información negativa sobre $A$ si se cumple que $P(A \mid B) \leq P(A)$.
        Demuestre que si $B$ con $0 < P(B) < 1$ aporta información negativa sobre $A$, entonces su complemento
        $B^c$ aporta información positiva, es decir, $P(A \mid B^c) \geq P(A)$.
    \end{mdframed}

    \begin{mdframed}[style=mdpinkbox,frametitle={Ejercicio 3}]
        Decimos que $A$ implica a $B$ probabilísticamente si se cumple que $P(B \mid A) = 1$. Demuestre
        que si dos eventos $A_1, A_2$ con la misma probabilidad inicial $P(A_1) = P(A_2)$ implican cada uno a $B$
        probabilísticamente, es decir, $P(B \mid A_i) = 1$ para $i = 1, 2$, entonces $A_1, A_2$ también tienen
        la misma probabilidad posterior $P(A_1 \mid B) = P(A_2 \mid B)$. Dé un ejemplo concreto en el que
        $P(B \mid A_1) = P(B \mid A_2) = 1$ pero $P(A_1 \mid B) \neq P(A_2 \mid B)$.
    \end{mdframed}

    \begin{mdframed}[style=mdpinkbox,frametitle={Ejercicio 4}]
        En lógica determinista, la afirmación “$A$ implica $B$” es equivalente a su contrapositiva “no $B$
        implica no $A$”. En este problema consideramos el equivalente probabilístico. Sean $A, B$ eventos
        con probabilidades distintas de $0$ y $1$.

        (a) Demuestre que $P(B \mid A) = 1$ es equivalente a $P(A^c \mid B^c) = 1$.

        (b) Dé un ejemplo donde lo anterior no se cumple si la igualdad es reemplazada por aproximadamente.
        Sugerencia: encuentre eventos independientes tales que $P(B \mid A)$ es muy cercano a $1$ pero
        $P(A^c \mid B^c)$ es muy cercano a $0$.
    \end{mdframed}

    \begin{mdframed}[style=mdpinkbox,frametitle={Ejercicio 5}]
        Demuestre que si $B$ es un evento tal que $P(B)(1 - P(B)) = 0$, entonces $B$ es independiente
        de cualquier otro evento $A$. En particular, si $P(B) = 1$, entonces para la probabilidad inicial
        y posterior de cualquier otro evento $A$ coinciden, esto es, $P(A \mid B) = P(A)$.
    \end{mdframed}

    \begin{mdframed}[style=mdpinkbox,frametitle={Ejercicio 6}]
        Demuestre que si $A_1, A_2, \dots, A_n$ son independientes, también lo son
        $A_1^c, A_2^c, \dots, A_n^c$. En particular, se cumple que
        \[
        P(A_1 \cup A_2 \cup \cdots \cup A_n) = 1 - \prod_{k=1}^{n} (1 - P(A_k)).
        \]
    \end{mdframed}

    \begin{mdframed}[style=mdpinkbox,frametitle={Ejercicio 7}]
        Alicia y Roberto comparten casa en una comunidad en la que habitan otras 100 personas.
        Si elegimos una persona al azar, la probabilidad de que tenga una amistad con Alicia es $\frac{1}{2}$,
        y lo mismo sucede con Roberto. Sea $C$ el conjunto de las otras 100 personas, $A$ el conjunto de amigos
        de Alicia y $B$ el conjunto de amigos de Roberto. Encuentre lo siguiente bajo la hipótesis de que
        todas las amistades son independientes.

        (a) $P(A = D)$ para algún $D \subset C$ dado.

        (b) $P(A \subset B)$.

        (c) $P(A \cup B = C)$.
    \end{mdframed}

    \begin{mdframed}[style=mdpinkbox,frametitle={Ejercicio 8}]
        En el problema de Monty Hall, el presentador da a elegir a un concursante una de tres puertas,
        detrás de las cuales hay dos cabras y un automóvil. Después de elegir, el presentador abre una
        puerta con cabra, quedando dos cerradas. Se puede conservar la elección o cambiar.

        Antes de cada concurso, Monty lanza una moneda con probabilidad $p \in (0,1)$ de obtener sol.
        El concursante conoce $p$ pero no el resultado. Si cae sol, Monty abre una puerta con cabra;
        si cae águila, abre una puerta al azar (posiblemente con automóvil).

        Si Monty ha abierto una puerta con una cabra, determine si es más conveniente conservar la elección
        inicial o cambiar de puerta.
    \end{mdframed}
\end{document}